\SourceLeaf{232}
casu classis ipsa \textit{primitiva} dicetur. Porro manifestum est, si forma aliqua $F$ determinantis $D$ derivata sit ex forma primitiva $f$ determinantis $\frac{D}{mm}$, classesque in quibus formae $F$, $f$ resp. contineantur, sint $K$, $k$, omnes formas e classe $K$ derivatas fore e classe primitiva $k$; quocirca classem $K$ ipsam \textit{ex classe primitiva $k$ derivatam} in hoc casu vocabimus.

Si $(a,b,c)$ est forma primitiva, neque vero $a$, $c$ simul pares (\textit{i. e.} si aut uterque impar aut saltem alteruter), facile intelligitur, non modo $a$, $b$, $c$, sed etiam $a$, $2b$, $c$ divisorem communem habere non posse, in quo casu forma $(a,b,c)$ dicitur \textit{proprie primitiva} sive simpliciter \textit{forma propria}. Si vero $(a,b,c)$ est forma primitiva, numeri $a$, $c$ autem ambo pares, patet, numeros $a$, $2b$, $c$ divisorem communem $2$ habere (qui simul erit maximus), vocabiturque $(a,b,c)$ \textit{forma improprie primitiva}, sive simpliciter \textit{forma impropria}\HistoricalNote{*)}{Hos terminos \textit{proprie} et \textit{improprie} ideo hic elegimus, quia alii magis idonei non occurrebant, quod admonemus, ne quis inter hanc significationem eamque qua inde ab art. 157 usi sumus, nexum occultum quaerat, qui nullus adest. Ceterum ambiguitas certe hinc non est metuenda.}. In hoc casu $b$ necessario erit impar (alioquin enim $(a,b,c)$ non esset forma primitiva); quare erit $bb\equiv1\ (\mathrm{mod.}\ 4)$ adeoque quoniam $ac$ per $4$ divisibilis, determinans $bb-ac\equiv1\ (\mathrm{mod.}\ 4)$. Formae impropriae itaque tantummodo pro determinante formae $4n+1$, si est positivus, vel formae $-(4n+3)$, si est negativus, locum habent. --- Ex art. 161 autem perspicuum est, si in classe aliqua data forma proprie primitiva inveniatur, omnes formas huius classis proprie primitivas esse; contra classem quae formam improprie primitivam implicet ex solis formis improprie primitivis constare. Quamobrem classis ipsa in casu priori \textit{proprie primitiva} seu simpliciter \textit{propria}; in posteriori \textit{improprie primitiva} seu \textit{impropria} appellabitur. Ita \textit{e. g.} inter classes positivas formarum determinantis $-235$ sex sunt propriae, puta quarum repraesentantes $(1,0,235)$, $(4,1,59)$, $(4,-1,59)$, $(5,0,47)$, $(13,5,20)$, $(13,-5,20)$, totidemque inter negativas; binae vero inter utrasque impropriae. --- Classes formarum determinantis $79$ (utpote numeri formae $4n+3$) omnes sunt propriae.

Si forma $(a,b,c)$ est derivata, et quidem e primitiva $\left(\frac{a}{m},\frac{b}{m},\frac{c}{m}\right)$, haec aut proprie primitiva aut improprie esse poterit. In casu priori $m$ erit divisor communis maximus etiam numerorum $a$, $2b$, $c$; in posteriori horum numerorum div. comm. max. erit $2m$. Hinc intelligitur distinctio inter \textit{formam e forma proprie primitiva derivatam} et \textit{formam ex improprie primitiva derivatam}; nec non (quoniam propter art. 161 omnes formae eiusdem classis hoc respectu perinde se habent)
\SourceLeaf{233}
inter \textit{classem derivatam e classe proprie primitiva} et \textit{classem ex improprie primitiva derivatam}.

Per has distinctiones fundamentum primum nacti sumus, cui distributionem omnium classium formarum determinantis dati in varios \textit{ordines} superstruere possumus. Classes duas, quarum repraesentantes sunt formae $(a,b,c)$, $(a',b',c')$ in \textit{eundem ordinem} coniiciemus, si \textit{tum} numeri $a$, $b$, $c$ eundem divisorem communem maximum habent ut $a'$, $b'$, $c'$, \textit{tum} $a$, $2b$, $c$ eundem ut $a'$, $2b'$, $c'$; si vero aut alterutra aut utraque harum conditionum locum non habet, classes ad \textit{ordines diversos} referentur. Hinc statim patet, omnes classes proprie primitivas unum ordinem constituere; omnes classes improprie primitivas, alium; si $mm$ est quadratum determinantem $D$ metiens, classes derivatae e classibus proprie primitivis determinantis $\frac{D}{mm}$ formabunt ordinem peculiarem, aliumque classes derivatae e classibus improprie primitivis determinantis $\frac{D}{mm}$ etc. Si forte $D$ per nullum quadratum (praeter $1$) divisibilis est, ordines classium derivatarum non aderunt adeoque aut unus tantum ordo dabitur (quando $D\equiv2$ vel $3$ secundum mod. $4$), puta ordo classium proprie primitivarum, aut duo (quando $D\equiv1\ (\mathrm{mod.}\ 4)$) scilicet O. classium proprie primitivarum et O. cl. impr. primitivarum. Per principia calculi combinationum haud difficile conditur regula sequens generalis: Si supponitur $D=D'2^{2\mu}a^{2\alpha}b^{2\beta}c^{2\gamma}\ldots.$ ita ut $D'$ nullum factorem quadraticum implicet, et $a$, $b$, $c$ etc. sint numeri primi impares diversi (ad quam formam quivis numerus redigi potest faciendo $\mu=0$, quando $D$ per $4$ non est divisibilis; et $\alpha$, $\beta$, $\gamma$ etc. omnes $=0$, sive quod eodem redit omittendo factores $a^{2\alpha}$, $b^{2\beta}$, $c^{2\gamma}$ etc., quando $D$ per nullum quadratum impar dividi potest): habebuntur aut ordines
\[
(\mu+1)(\alpha+1)(\beta+1)(\gamma+1)\ldots
\]
nempe quando $D'\equiv2$ vel $3\ (\mathrm{mod.}\ 4)$; aut ordines
\[
(\mu+2)(\alpha+1)(\beta+1)(\gamma+1)\ldots
\]
quando $D'\equiv1\ (\mathrm{mod.}\ 4)$. Sed demonstrationem huius regulae supprimimus, quoniam neque difficilis neque hic adeo necessaria est.

\textit{Ex.} 1. Pro $D=45=5.3^2$ habentur sex classes, quarum repraesentantes $(1,0,-45)$, $(-1.\,0,45)$, $(2,1,-22)$, $(-2,1,22)$, $(3.\,0,-15)$, $(6,3,-6)$. Hae distribuuntur in quatuor ordines, scilicet O. I comprehendet duas classes proprias quarum repr. $(1,0,-45)$, $(-1,0,45)$; O. II continebit duas classes im%
\SourceLeaf{234}%
proprias, quarum repr. $(2,1,-22)$, $(-2,1,22)$; O. III continebit unam classem derivatam e propria determinantis $5$, puta cuius repr. $(3,0,-15)$; O. IV constabit ex una classe derivata ex impropria det. $5$, puta cuius repr. $(6,3,-6)$.

\textit{Ex.} 2. Classes positivae determinantis $-99=-11.3^2$ inter quatuor ordines distribuentur: O. I complectetur classes proprie primitivas sequentes\HistoricalNote{*)}{Adhibendo brevitatis caussa formas repraesentantes pro classibus ipsis quarum vice funguntur.}: $(1,0,99)$, $(4,1,25)$. $(4,-1,25)$, $(5,1,20)$, $(5,-1,20)$, $(9,0,11)$; O. II continebit classes improprias $(2,1,50)$, $(10,1,10)$; O. III classes derivatas e propriis determinantis $-11$, $(3,0,33)$, $(9,3,12)$, $(9,-3,12)$; O. IV classem unicam derivatam ex impropria det. $-11$, $(6,3,18)$. --- Classes negativae huius determinantis prorsus eodem modo in ordines distribui poterunt.

Observamus, \textit{classes oppositas semper ad eundem ordinem referri}, cuius theorematis ratio nullo negotio perspicitur.

\Needspace{8\baselineskip}
\Article{227}
Ex his diversis ordinibus imprimis ordo classium proprie primitivarum maximam attentionem meretur. Nam singulae classes derivatae a certis classibus primitivis (determinantis minoris) originem trahunt, ex quarum consideratione ea quae ad illas spectant plerumque sponte sequuntur. Infra autem docebimus, quamlibet classem improprie primitivam simili modo quasi associatam esse aut unicae classi proprie primitivae aut tribus (eiusdem determinantis). Porro pro determinantibus negativis classes negativas praeterire licebit, quippe quibus singulis certae classes positivae semper respondent. Ut itaque naturam classium proprie primitivarum profundius penetremus, ante omnia differentiam certam essentialem explicabimus, secundum quam totus ordo classium propriarum in plura \textit{genera} subdividi potest. Quoniam hoc argumentum gravissimum hactenus nondum attigimus, res ab integro nobis erit repetenda.

\Needspace{12\baselineskip}
\begin{center}
\textit{Ordinum partitio in genera.}
\end{center}

\Needspace{8\baselineskip}
\Article{228}
\textsc{Theorema.} \textit{Per formam quamcunque proprie primitivam $F$ repraesentari possunt infinite multi numeri per numerum primum quemcunque datum $p$ non divisibiles.}

\textit{Dem.} Si forma $F=axx+2bxy+cyy$, manifestum est, $p$ omnes tres numeros $a$, $2b$, $c$ simul metiri non posse. Iam quando $a$ per $p$ non est divisi%
\SourceLeaf{235}%
bilis, patet. si pro $x$ assumatur numerus quicunque per $p$ non divisibilis, pro $y$ vero numerus per $p$ divisibilis, valorem formae $F$ fieri non divisibilem per $p$; quando $c$ per $p$ non est divisibilis, idem obtinetur tribuendo ipsi $x$ valorem divisibilem ipsique $y$ valorem non divisibilem; denique quando tum $a$ tum $c$ per $p$ sunt divisibiles, adeoque $2b$ non divisibilis, forma $F$ valorem per $p$ non divisibilem induet tribuendo tum ipsi $x$ tum ipsi $y$ valores quoscunque per $p$ non divisibiles. \textit{Q. E. D.}

Manifestum est, theorema etiam pro formis \textit{improprie primitivis} locum habere, si modo non fuerit $p=2$.

Quoniam plures huiusmodi conditiones simul consistere possunt, ut idem numerus per quosdam numeros primos datos divisibilis sit, per alios non divisibilis (v. art. 32): facile perspicitur, numeros $x$, $y$ infinite multis modis ita determinari posse, ut forma primitiva $axx+2bxy+cyy$ valorem per quotcunque numeros primos datos non divisibilem adipiscatur, a quibus unice excludendus est $2$, quoties forma est improprie primitiva. Hinc patet, theorema generalius ita proponi posse: \textit{Per formam quamcunque primitivam repraesentari possunt infinite multi numeri, qui ad numerum quemcunque datum (imparem, quando forma est improprie primitiva) sint primi.}

\Needspace{8\baselineskip}
\Article{229}
\textsc{Theorema.} \textit{Sit $F$ forma primitiva determinantis $D$, $p$ numerus primus ipsum $D$ metiens: tum numeri per $p$ non divisibiles qui per formam $F$ repraesentari possunt, in eo convenient, ut vel omnes sint residua quadratica ipsius $p$, vel omnes non-residua.}

\textit{Dem.} Sit $F=(a,b,c)$; $m$, $m'$ duo numeri quicunque per $p$ non divisibiles qui per formam $F$ repraesentari possunt, scilicet
\[
m=agg+2bgh+chh,\qquad m'=ag'g'+2bg'h'+ch'h'
\]
Tum erit
\[
mm'=(agg'+b(gh'+hg')+chh')^2-D(gh'-hg')^2
\]
quare $mm'$ quadrato congruus erit secundum modulum $D$, adeoque etiam secundum $p$, \textit{i. e.} $mm'$ erit residuum quadraticum ipsius $p$. Hinc sequitur, aut utrumque $m$, $m'$ esse residuum quadraticum ipsius $p$, aut utrumque non-residuum. \textit{Q. E. D.}

\SourceLeaf{236}
Simili modo probatur, quando determinans $D$ per $4$ sit divisibilis, omnes numeros impares per $F$ repraesentabiles vel esse $\equiv1$. vel omnes $\equiv3\ (\mathrm{mod.}\ 4)$. Scilicet productum e duobus numeris talibus in hoc casu semper erit residuum quadr. ipsius $4$, adeoque $\equiv1\ (\mathrm{mod.}\ 4)$; quare vel uterque erit $\equiv1$, vel uterque $\equiv3$.

Denique quando $D$ per $8$ est divisibilis, productum e duobus numeris quibuscunque imparibus, qui per $F$ repraesentari possunt, erit R. Q. ipsius $8$ et proin $\equiv1\ (\mathrm{mod.}\ 8)$. Quare in hoc casu omnes numeri impares per $F$ repraesentabiles vel erunt $\equiv1$, vel omnes $\equiv3$, vel omnes $\equiv5$, vel omnes $\equiv7\ (\mathrm{mod.}\ 8)$.

Ita \textit{e. g.} quum per formam $(10,3,17)$ repraesentari possit numerus $10$ qui est N. R. ipsius $7$: omnes numeri per $7$ non divisibiles, qui per formam illam repraesentari possunt, non-residua ipsius $7$ erunt. --- Quum $-3$ per formam $(-3,1,49)$ repraesentabilis et sec. mod. $4$ sit $\equiv1$, omnes numeri impares per formam hanc repraesentabiles perinde se habebunt.

Ceterum, si ad propositum praesens necessarium esset, facile demonstrare possemus, numeros per formam $F$ repraesentabiles ad nullum numerum primum qui ipsum $D$ non metiatur, talem relationem fixam habere, sed promiscue tum residua tum non-residua numeri cuiusvis primi ipsum $D$ non metientis per formam $F$ repraesentari posse. Contra respectu numerorum $4$ et $8$ analogon quoddam etiam in aliis casibus locum habet, quos praeterire non possumus.

I. \textit{Quando determinans $D$ formae primitivae $F$ est $\equiv3\ (\mathit{mod.}\ 4)$: omnes numeri impares, per formam $F$ repraesentabiles, erunt vel $\equiv1$, vel omnes $\equiv3\ (\mathit{mod.}\ 4)$.} Si enim $m$, $m'$ sunt duo numeri per $F$ repraesentabiles, productum $mm'$ eodem modo ut supra sub formam $pp-Dqq$ redigi poterit. Quando itaque uterque $m$, $m'$ est impar, necessario alter numerorum $p$, $q$ par erit, alter impar adeoque alterum quadratorum $pp$, $qq$, $\equiv0$, alterum $\equiv1\ (\mathrm{mod.}\ 4)$. Unde facile deducitur, $pp-Dqq$ certo esse $\equiv1\ (\mathrm{mod.}\ 4)$, adeoque aut utrumque $m$, $m'$, $\equiv1$, aut utrumque $\equiv3\ (\mathrm{mod.}\ 4)$. Ita \textit{e. g.} per formam $(10,3,17)$ alii numeri impares quam qui sunt formae $4n+1$ repraesentari nequeunt.

II. \textit{Quando determinans $D$ formae primitivae $F$ est $\equiv2\ (\mathit{mod.}\ 8)$: omnes numeri impares, per formam $F$ repraesentabiles, erunt vel partim $\equiv1$ partim $\equiv7$, vel partim $\equiv3$ partim $\equiv5\ (\mathit{mod.}\ 8)$.} Ponamus enim $m$, $m'$ esse duos numeros
\SourceLeaf{237}
impares per $F$ repraesentabiles, quorum igitur productum $mm'$ sub formam $pp-Dqq$ redigi poterit. Quando ergo uterque $m$, $m'$ est impar, necessario $p$ impar esse debebit (quia $D$ par), adeoque $pp\equiv1\ (\mathrm{mod.}\ 8)$; $qq$ vero erit vel $\equiv0$ vel $\equiv1$ vel $\equiv4$, et proin $Dqq$ vel $\equiv0$ vel $\equiv2$. Hinc $mm'=pp-Dqq$ fit vel $\equiv1$ vel $\equiv7\ (\mathrm{mod.}\ 8)$; si itaque $m$ est vel $\equiv1$ vel $\equiv7$, etiam $m'$ erit vel $\equiv1$ vel $\equiv7$; si vero $m$ est vel $\equiv3$ vel $\equiv5$, etiam $m'$ erit vel $\equiv3$ vel $\equiv5$. \textit{E. g.} omnes numeri impares per formam $(3,1,5)$ repraesentabiles sunt aut $\equiv3$ aut $\equiv5\ (\mathrm{mod.}\ 8)$, nullique numeri formae $8n+1$ aut $8n+7$ per formam illam repraesentari possunt.

III. \textit{Quando determinans $D$ formae primitivae $F$ est $\equiv6\ (\mathit{mod.}\ 8)$: per formam hanc repraesentari possunt numeri impares vel tales tantum qui sunt $\equiv1$ et $\equiv3\ (\mathit{mod.}\ 8)$, vel tales tantum qui sunt $\equiv5$ et $\equiv7\ (\mathit{mod.}\ 8)$.} Demonstrationem praecedenti (in II) omnino similem quisque nullo negotio evolvere poterit. --- Ita \textit{e. g.} per formam $(5,1,7)$ unice tales numeri impares possunt repraesentari qui sunt aut $\equiv5$ aut $\equiv7\ (\mathrm{mod.}\ 8)$.

\Needspace{8\baselineskip}
\Article{230}
Omnes igitur numeri qui per formam primitivam datam $F$ determinantis $D$ repraesentari possunt, relationem fixam habebunt ad singulos divisores primos ipsius $D$ (per quos quidem ipsi non sunt divisibiles), numeri impares vero qui per $F$ possunt repraesentari, in quibusdam casibus etiam ad numeros $4$ et $8$ relationem fixam habebunt, scilicet ad $4$, quoties $D$ aut $\equiv0$ aut $\equiv3\ (\mathrm{mod.}\ 4)$, et ad $8$, quoties $D$ aut $\equiv0$, aut $\equiv2$ aut $\equiv6\ (\mathrm{mod.}\ 8)$\HistoricalNote{*)}{Pro determinantibus per $8$ divisibilibus relatio ad numerum $4$ negligi potest, quoniam in hoc casu sub relatione ad $8$ iam est contenta.}. Talem relationem ad singulos hos numeros, \textit{characterem} seu \textit{characterem particularem} formae $F$ vocabimus sequentique modo exprimemus: Quando sola residua quadratica numeri primi $p$ per formam $F$ repraesentari possunt, tribuemus ipsi characterem $Rp$, in casu opposito characterem $Np$; similiter scribemus $1,4$, quando alii numeri impares per formam $F$ repraesentari nequeunt nisi qui sunt $\equiv1\ (\mathrm{mod.}\ 4)$, unde statim liquet quales characteres exprimantur per signa $3,4$; $1,8$; $3,8$; $5,8$; $7,8$. Denique formis per quas numeri impares tales soli repraesentari possunt qui sec.
\SourceLeaf{238}
mod. $8$ sunt vel $\equiv1$ vel $\equiv7$, tribuemus characterem $1\ \mathit{et}\ 7,8$; ex quo significatio characterum $3\ \mathit{et}\ 5,8$; $1\ \mathit{et}\ 3,8$; $5\ \mathit{et}\ 7,8$ sponte sequitur.

Characteres singuli formae primitivae datae $(a,b,c)$ determinantis $D$ semper ex uno saltem numerorum $a$, $c$ (qui manifesto per formam illam ambo sunt repraesentabiles) cognosci possunt. Nam quoties $p$ est divisor primus ipsius $D$, certe unus numerorum $a$, $c$ per $p$ non erit divisibilis; si enim uterque per $p$ divisibilis esset, $p$ etiam ipsum $bb(=D+ac)$ metiretur. et proin etiam ipsum $b$, \textit{i. e.} forma $(a,b,c)$ non esset primitiva. Simili modo in iis casibus, ubi forma $(a,b,c)$ ad numerum $4$ vel $8$ relationem fixam habet, certo ad minimum unus numerorum $a$, $c$ impar erit, ex quo igitur relatio illa deprehendi poterit. Ita \textit{e. g.} character formae $(7,0,23)$ respectu numeri $23$ e numero $7$ concluditur $N23$, eiusdem formae character respectu numeri $7$ habetur ex numero $23$ puta $R7$; denique character huius formae respectu numeri $4$, puta $3,4$, vel e numero $7$ vel e numero $23$ colligi potest.

Quoniam omnes numeri qui per formam aliquam $F$ in classe $K$ contentam repraesentari possunt, etiam per quamlibet aliam formam huius classis sunt repraesentabiles: manifesto singuli characteres formae $F$ omnibus reliquis formis huius classis quoque competent, quapropter illos tanquam characteres totius classis considerare licebit. Singuli itaque characteres classis cuiuslibet primitivae datae ex ipsius forma repraesentante cognoscuntur. Classes oppositae semper characteres omnes eosdem habebunt.

\Needspace{8\baselineskip}
\Article{231}
Complexus \textit{omnium} characterum particularium formae vel classis datae constituet characterem integrum huius formae vel classis. Ita \textit{e. g.} character integer formae $(10,3,17)$, vel totius classis quam repraesentat erit $1,4$; $N7$; $N23$. Simili modo character integer formae $(7,1,-17)$ erit $7,8$; $R3$; $N5$, nam character particularis $3,4$ in hoc casu omittitur quia in charactere $7,8$ iam est contentus. --- Ex hoc fonte petimus subdivisionem totius ordinis classium proprie primitivarum (positivarum quando det. est negativus) determinantis dati in plura \textit{genera} diversa, referendo omnes classes, quae eundem characterem integrum habent, ad genus idem; quarumque characteres integri diversi sunt, ad genera diversa. Singulis vero generibus eos characteres integros tribuemus, quos classes sub ipsis contentae habent. Ita \textit{e. g.} pro determinante $-161$ habentur sedecim
\SourceLeaf{239}
classes positivae proprie primitivae. quae sequenti modo in quatuor genera distribuuntur:
\begin{center}
\begin{tabular}{c|l}
Character & \multicolumn{1}{c}{Classium formae repraesentantes} \\
$1,4;\ R7;\ R23$ & $(1,0,161)$, $(2,1,81)$, $(9,1,18)$, $(9,-1,18)$ \\
$1,4;\ N7;\ N23$ & $(5,2,33)$, $(5,-2,33)$, $(10,3,17)$, $(10,-3,17)$ \\
$3,4;\ R7;\ N23$ & $(7,0,23)$, $(11,2,15)$, $(11,-2,15)$, $(14,7,15)$ \\
$3,4;\ N7;\ R23$ & $(3,1,54)$, $(3,-1,54)$, $(6,1,27)$, $(6,-1,27)$
\end{tabular}
\end{center}

De multitudine characterum integrorum diversorum, qui quidem a priori sunt possibiles, teneantur sequentia.

I. Quando determinans $D$ per $8$ est divisibilis, respectu numeri $8$ quatuor characteres particulares diversi sunt possibiles; numerus $4$ nullum characterem peculiarem suppeditat (annot. ad art. praec.). Praeterea respectu singulorum divisorum primorum imparium ipsius $D$ bini characteres dantur; quare si illorum multitudo est $m$, dabuntur omnino $2^{m+2}$ characteres integri diversi (statuendo $m=0$, quoties $D$ est potestas binaria).

II. Quando det. $D$ per $8$ non est divisibilis, sed tamen per $4$, insuperque per $m$ numeros primos impares: omnino habebuntur $2^{m+1}$ characteres integri diversi.

III. Quando det. $D$ est par neque vero per $4$ divisibilis, erit vel $\equiv2\ (\mathrm{mod.}\ 8)$ vel $\equiv6$. In casu priori dabuntur duo characteres particulares respectu numeri $8$ puta $1\ \mathit{et}\ 7,8$, atque $3\ \mathit{et}\ 5,8$; in casu posteriori totidem. Posita igitur multitudine divisorum primorum imparium ipsius $D$, $=m$: habebuntur omnino $2^{m+1}$ characteres integri diversi.

IV. Quando $D$ est impar, erit vel $\equiv1$ vel $\equiv3\ (\mathrm{mod.}\ 4)$. In casu posteriori respectu numeri $4$ duo characteres diversi dantur, qualis relatio in casu priori in characterem integrum non ingreditur. Quare designante $m$ idem ut ante, in casu priori dabuntur $2^m$, in posteriori $2^{m+1}$ characteres integri diversi.

Probe vero notandum est, hinc neutiquam sequi, totidem genera revera
\SourceLeaf{240}
dari quot characteres diversi a priori sint possibiles. In exemplo quidem nostro horum semissi tantum revera classes sive genera respondent, nullaeque classes positivae dantur, quibus characteres $1,4;\ R7;\ N23$ vel $1,4;\ N7;\ R23$; vel $3,4;\ R7;\ R23$ vel $3,4;\ N7;\ N23$ competant. De quo argumento gravissimo infra fusius agetur.

Formae $(1,0,-D)$, quae haud dubie inter omnes formas determinantis $D$ pro simplicissima habenda est, nomen \textit{formae principalis} abhinc tribuemus; classem totam in qua illa reperitur, \textit{classem principalem} vocabimus; denique genus totum in quo classis principalis contenta est, \textit{genus principale} dicetur. Probe itaque distinguendae sunt forma principalis, forma e classe principali et forma e genere principali; nec non classis principalis et classis e genere principali. His denominationibus semper utemur, etiamsi forte pro determinante aliquo aliae classes praeter principalem, vel alia genera praeter genus principale non dentur, uti \textit{e. g.} evenit plerumque, quando $D$ est numerus primus positivus formae $4n+1$.

\Needspace{8\baselineskip}
\Article{232}
Quamquam ea quae de formarum characteribus explicata sunt proxime eum in finem sunt allata, ut subdivisio ordinis \textit{positivi proprie primitivi} inde petatur: tamen nihil impedit quominus eadem etiam ad formas classesque negativas aut ad improprie primitivas applicentur, atque tum ordo improprie primitivus positivus, tum ordo proprie primitivus negativus, tum ordo improprie primitivus negativus ex eodem principio in genera subdividantur. Ita postquam \textit{e. g.} ordo proprie primitivus formarum determinantis $145$ in duo genera sequentia subdivisus est
\begin{center}
\begin{tabular}{l|l}
$R5,\ R29$ & $(1,0,-145)$, $(5,0,-29)$ \\
$N5,\ N29$ & $(3,1,-48)$, $(3,-1,-48)$
\end{tabular}
\end{center}
etiam ordo improprie primitivus perinde in duo genera subdividi potest:
\begin{center}
\begin{tabular}{l|l}
$R5,\ R29$ & $(4,1,-36)$, $(4,-1,-36)$ \\
$N5,\ N29$ & $(2,1,-72)$, $(10,5,-12)$
\end{tabular}
\end{center}
vel, sicuti classes positivae formarum determinantis $-129$ in quatuor genera distribuuntur:
\SourceLeaf{241}
\begin{center}
\begin{tabular}{l|l}
$1,4;\ R3;\ R43$ & $(1,0,129)$, $(10,1,13)$, $(10,-1,13)$ \\
$1,4;\ N3;\ N43$ & $(2,1,65)$, $(5,1,26)$, $(5,-1.\,26)$ \\
$3,4;\ R3;\ N43$ & $(3,0,43)$, $(7,2,19)$, $(7,-2,19)$ \\
$3,4;\ N3;\ R43$ & $(6,3,23)$, $(11,5,14)$, $(11,-5,14)$
\end{tabular}
\end{center}
etiam classes negativae in quatuor ordines discedunt
\begin{center}
\begin{tabular}{l|l}
$3,4;\ N3;\ N43$ & $(-1,0,-129)$, $(-10,1,-13)$, $(-10,-1,-13)$ \\
$3,4;\ R3;\ R43$ & $(-2,1,-65)$, $(-5,1,-26)$, $(-5,-1,-26)$ \\
$1,4;\ N3;\ R43$ & $(-3,0,-43)$, $(-7,2,-19)$, $(-7,-2,-19)$ \\
$1,4;\ R3;\ N43$ & $(-6,3,-23)$, $(-11,5,-14)$, $(-11,-5,-14)$
\end{tabular}
\end{center}
Attamen quum systema classium negativarum systemati positivarum semper tam simile evadat, plerumque superfluum videbitur illud seorsim construere. Ordinem improprie primitivum autem ad proprie primitivum reducere infra docebimus.

Tandem quod attinet ad ordines derivatos: pro harum subdivisione regulae novae non sunt necessariae. Quum enim quivis ordo derivatus ex aliquo ordine primitivo (determinantis minoris) originem trahat, illiusque classes singulae ad singulas huius sponte referantur: manifesto subdivisio ordinis derivati e subdivisione ordinis primitivi peti poterit.

\Needspace{8\baselineskip}
\Article{233}
Si forma (primitiva) $F=(a,b,c)$ ita est comparata, ut inveniri possint duo numeri $g$, $h$ tales ut fiat $gg\equiv a$, $gh\equiv b$, $hh\equiv c$ secundum modulum datum $m$, dicemus formam illam esse residuum quadraticum numeri $m$ atque $gx+hy$ valorem expressionis $\surd(axx+2bxy+cyy)(\mathrm{mod.}\ m)$, sive brevius $(g,h)$ valorem expr. $\surd(a,b,c)$ vel $\surd F(\mathrm{mod.}\ m)$. Generalius, si multiplicator $M$, ad modulum $m$ primus, eius est indolis ut fieri possit
\[
gg\equiv aM,\quad gh\equiv bM,\quad hh\equiv cM\quad (\mathrm{mod.}\ m)
\]
dicemus $M\times(a,b,c)$ sive $MF$ esse res. quadr. ipsius $m$, atque $(g,h)$ valorem expressionis $\surd M(a,b,c)$ vel $\surd MF(\mathrm{mod.}\ m)$. Ita \textit{e. g.} forma $(3,1,54)$ est res. quadr. ipsius $23$ atque $(7,10)$ valor expr. $\surd(3,1,54)(\mathrm{mod.}\ 23)$; similiter $(2,-4)$ valor expr. $\surd5(10,3,17)(\mathrm{mod.}\ 23)$. Usus harum definitionum infra ostendetur: hic notentur propositiones sequentes:
